% Epilogue of The Replica Rosetta Stone (written by the coordinating author from the chapters'
% final status boards; figures/scripts/16-epilogue-boards.py reads the boards and counts them).
\specialchapter{Epilogue}{The stone, read across}{open}
\rstfloatprefix{E}
\begin{refsection}[chapters/16-epilogue.bib]

The prologue derived one number, the capacity $2$ of the perceptron, in four languages, and found that the physicist's replicas and the probabilist's comparison inequality, set side by side, do the same work in the same order. Fifteen chapters later the same exercise has been carried out on fifteen predictions, from the free entropy of the simplest disordered systems to the plateaus of stochastic gradient descent. Each chapter closed its account with a status board that says, in one word per prediction, how much of it is known. Read one at a time, the boards record theorems and conjectures. Read side by side, they show a pattern: which physical steps the proofs justify, which steps they replace, and which steps they have not been able to touch. This epilogue sets the boards side by side, reads the stone across the chapters, and collects what is false, what rests on unrefereed preprints, and what is open.

\section*{The fifteen predictions}
\addcontentsline{toc}{section}{The fifteen predictions}

\Cref{tab:epilogue:predictions,tab:epilogue:predictionsb} list the main prediction of the prologue and of each chapter, the derivation that produced it, the route by which it was proved, and the status of the stone at the foot of the opening page. The precise statements, with their hypotheses, are in the chapters.

\begin{table}[tp]
\begin{fullwidth}
\footnotesize
\renewcommand{\arraystretch}{1.2}
\begin{tabularx}{\linewidth}{@{}>{\raggedright\arraybackslash}p{6mm}>{\hsize=1.05\hsize\raggedright\arraybackslash}X>{\hsize=0.95\hsize\raggedright\arraybackslash}X>{\hsize=1.0\hsize\raggedright\arraybackslash}X>{\raggedright\arraybackslash}p{17mm}@{}}
\toprule
& \textsc{\color{stone}the prediction} & \textsc{\color{physics}derived by} & \textsc{\color{proof}proved by} & \textsc{\color{stone}stone}\\
\midrule
P & the capacity of the spherical perceptron: $2$ at $\kappa=0$, and $\alpha_c(\kappa)$ for $\kappa\ge0$ & Gardner's replicas \parencite*{Gardner1988}, after Cover's count at $\kappa=0$ & Cover's count at $\kappa=0$; Gordon's comparison \parencite{Stojnic2013a}; the volume \parencite{ShcherbinaTirozzi2003}; conic geometry \parencite{AmelunxenEtAl2014} & \statusstone{proved}\enspace proved\\
\multicolumn{5}{@{}l}{\textsc{\color{stone}part i. two languages}}\\[0.2ex]
1 & the free entropy self-averages and is a saddle point & Hubbard--Stratonovich and the replica trick & Gaussian concentration; Laplace's method on types & \statusstone{proved}\enspace proved\\
2 & the REM freezes at $\beta_c=\sqrt{2\log2}$, with one-step breaking below & \textcite{Derrida1980,Derrida1981}; \textcite{GrossMezard1984} & first and second moments; Poisson--Dirichlet weights \parencite{Ruelle1987,BovierKurkova2004} & \statusstone{proved}\enspace proved\\
3 & the SK free entropy: replica symmetric above the de~Almeida--Thouless line, Parisi's formula everywhere & \textcite{SherringtonKirkpatrick1975}; \textcite{deAlmeidaThouless1978}; \textcite{Parisi1979b} & interpolation: \textcite{GuerraToninelli2002}, \textcite{Guerra2003}, \textcite{Talagrand2006}; the line with a field in 2026 preprints & \statusstone{proved}\enspace proved\\
\multicolumn{5}{@{}l}{\textsc{\color{stone}part ii. inference on the nishimori line}}\\[0.2ex]
4 & the mutual information of $\Z_2$ synchronization is replica symmetric & the replica formula on the Nishimori line \parencite{LesieurKrzakalaZdeborova2017} & interpolation and overlap concentration \parencite{DeshpandeAbbeMontanari2016,LelargeMiolane2019,BarbierMacris2019} & \statusstone{proved}\enspace proved\\
5 & PCA sees the spike above $\lambda=1$: eigenvalue $\lambda+1/\lambda$, overlap $1-1/\lambda^2$ & an eigenvalue can split off when the entries have a mean \parencite{JonesKosterlitzThouless1978}; replicas for covariances \parencite{HoyleRattray2004} & large means \parencite{FurediKomlos1981}; at the threshold \parencite{BaikBenArousPeche2005,FeralPeche2007,BenaychGeorgesNadakuditi2011}, here by the resolvent & \statusstone{proved}\enspace proved\\
6 & AMP's iterates are Gaussian with the variance of state evolution & the cavity and TAP \parencite{ThoulessAndersonPalmer1977,DonohoMalekiMontanari2009} & Gaussian conditioning \parencite{Bolthausen2014,BayatiMontanari2011}, for a fixed number of iterations & \statusstone{proved}\enspace proved\\
7 & the Bayes error of the teacher--student perceptron; perfect learning of the Ising teacher at $\alpha_{\mathrm{IT}}=1.2447$ & \textcite{GardnerDerrida1989b}; \textcite{Gyorgyi1990}; \textcite{SompolinskyTishbySeung1990} & adaptive interpolation \parencite{BarbierEtAl2019}, under hypotheses (h1)--(h5) & \statusstone{partial}\enspace in part\\
\bottomrule
\end{tabularx}
\end{fullwidth}
\caption{The main prediction of the prologue and of each chapter of Parts~I and~II, who derived it and how, who proved it and how, and the status of the stone at the foot of the opening page (``in part'': proved under extra hypotheses or in part of the range).}\label{tab:epilogue:predictions}
\end{table}

\begin{table}[tp]
\begin{fullwidth}
\footnotesize
\renewcommand{\arraystretch}{1.2}
\begin{tabularx}{\linewidth}{@{}>{\raggedright\arraybackslash}p{6mm}>{\hsize=1.05\hsize\raggedright\arraybackslash}X>{\hsize=0.95\hsize\raggedright\arraybackslash}X>{\hsize=1.0\hsize\raggedright\arraybackslash}X>{\raggedright\arraybackslash}p{17mm}@{}}
\toprule
& \textsc{\color{stone}the prediction} & \textsc{\color{physics}derived by} & \textsc{\color{proof}proved by} & \textsc{\color{stone}stone}\\
\midrule
\multicolumn{5}{@{}l}{\textsc{\color{stone}part iii. comparison and convexity}}\\[0.2ex]
8 & Gardner's capacity for every margin & Gardner's volume \parencite*{Gardner1988} & Gordon's min--max and the CGMT, for $\kappa\ge0$ \parencite{Gordon1985,Stojnic2013a,ThrampoulidisAbbasiHassibi2018} & \statusstone{partial}\enspace in part\\
9 & the ridge risk and its peak at $\alpha=1$ & replicas at zero temperature; random matrices \parencite{KroghHertz1992} & the Marchenko--Pastur resolvent \parencite{DobribanWager2018,HastieEtAl2022}; the CGMT & \statusstone{proved}\enspace proved\\
10 & the LASSO risk as a scalar fixed point; the Donoho--Tanner threshold & replicas \parencite{KabashimaEtAl2009,RanganFletcherGoyal2012} & state evolution \parencite{BayatiMontanari2012}; the CGMT \parencite{ThrampoulidisAbbasiHassibi2018}; conic geometry \parencite{AmelunxenEtAl2014} & \statusstone{proved}\enspace proved\\
11 & the logistic MLE exists exactly above $\alpha_{\mathrm{MLE}}(r)$; the max-margin error & statisticians, without physics \parencite{CandesSur2020}; the error: \textcite{OpperEtAl1990}; Gardner's volume with a teacher in \cref{ch:classification} & the kinematic formula \parencite{CandesSur2020}; the CGMT \parencite{MontanariEtAl2025} & \statusstone{proved}\enspace proved\\
12 & kernel learning curves through an effective ridge; Gaussian equivalence & replicas and cavities \parencite{BordelonEtAl2020,CanatarEtAl2021,GeraceEtAl2020} & deterministic equivalents \parencite{ChengMontanari2024}; universality \parencite{MeiMontanari2022,HuLu2023} & \statusstone{partial}\enspace in part\\
\multicolumn{5}{@{}l}{\textsc{\color{stone}part iv. landscapes and dynamics}}\\[0.2ex]
13 & the complexity of the spherical $p$-spin; threshold $-E_\infty$ above the ground state $-E_0$ & TAP complexity \parencite{CrisantiSommers1995,CavagnaGiardinaParisi1998} & Kac--Rice and the GOE \parencite{AuffingerBenArousCerny2013}; the second moment \parencite{Subag2017a} & \statusstone{partial}\enspace in part\\
14 & closed equations for correlation and response; dynamics stops at the threshold & the generating functional \parencite{CrisantiHornerSommers1993,CugliandoloKurchan1993} & large deviations \parencite{BenArousDemboGuionnet2006}; \textcite{Sellke2024}; conditioning \parencite{GerbelotEtAl2024} & \statusstone{partial}\enspace in part\\
15 & the Saad--Solla equations and the plateau of online SGD & order-parameter equations \parencite{SaadSolla1995a,BiehlSchwarze1995} & martingales on fixed horizons \parencite{GoldtEtAl2020a,VeigaEtAl2022} & \statusstone{proved}\enspace proved\\
\bottomrule
\end{tabularx}
\end{fullwidth}
\caption{The main prediction of each chapter of Parts~III and~IV, as in \cref{tab:epilogue:predictions}.}\label{tab:epilogue:predictionsb}
\end{table}

Of the sixteen stones of the prologue and the chapters, eleven read \emph{proved} and five read \emph{in part}. None reads \emph{conjectured} or \emph{false}: every main prediction of the book is a theorem in at least part of its range. This reflects the plan, which chose predictions whose proofs a chapter could present, not the state of the field. The boards are less selective. They list $149$ predictions in all, the statements that the chapters derived or discussed, and \cref{fig:epilogue:boards} counts their statuses. Seventy-one are proved, twenty-four more are proved under extra hypotheses, nine are proved only in unrefereed preprints, fifteen are conjectured, eight are heuristic, nineteen are known to be false, and three are open.

\begin{figure}[t]
\centering
\begin{tikzpicture}
\begin{axis}[rplot,xbar stacked,bar width=6pt,height=0.78\linewidth,
  xmin=0,xmax=13.5,ymin=-0.6,ymax=15.6,y dir=reverse,
  ytick={0,...,15},yticklabels={P,1,2,3,4,5,6,7,8,9,10,11,12,13,14,15},
  xlabel={rows on the chapter's status board},ylabel={chapter},
  legend style={at={(0.99,0.99)},anchor=north east,font=\scriptsize},legend cell align=left]
\addplot[fill=proof,draw=proof] table[x=proved,y=num] {figures/data/16-epilogue-boards.dat};
\addplot[fill=proof!35,draw=proof] table[x=under,y=num] {figures/data/16-epilogue-boards.dat};
\addplot[fill=proofwash,draw=proof,densely dashed] table[x=preprint,y=num] {figures/data/16-epilogue-boards.dat};
\addplot[fill=physics!30,draw=physics] table[x=conj,y=num] {figures/data/16-epilogue-boards.dat};
\addplot[fill=stone!25,draw=stone] table[x=heur,y=num] {figures/data/16-epilogue-boards.dat};
\addplot[fill=physics,draw=physics] table[x=false,y=num] {figures/data/16-epilogue-boards.dat};
\addplot[fill=white,draw=stone] table[x=open,y=num] {figures/data/16-epilogue-boards.dat};
\legend{proved,proved under,preprint,conjectured,heuristic,known false,open}
\end{axis}
\end{tikzpicture}
\caption{The statuses on the boards of the prologue (P) and the fifteen chapters, counted by \texttt{figures/scripts/16-epilogue-boards.py}, which reads every \texttt{\textbackslash statusrow} of the chapters' sources. Lapis: proved, in a refereed source or in the book (filled), under extra hypotheses (light), or only in an unrefereed preprint (dashed). Carnelian: conjectured (light) or known to be false (filled). Granite: heuristic (light) or open (white).}\label{fig:epilogue:boards}
\end{figure}

Two features of the counts stand out. The first is how much is proved: almost two thirds of the rows, $95$ of $149$, are theorems, with or without extra hypotheses. Physicists derived most of these first, and in most cases a decade or more passed before the proof: Gardner's capacity of 1988 waited until 2003 for $\kappa>0$ (for patterns with entries $\pm1$), Parisi's formula of 1979 until 2006, the Saad--Solla equations of 1995 until 2020, with gaps that \textcite{VeigaEtAl2022} closed, and the threshold energy that \textcite{CugliandoloKurchan1993} found in the dynamics until \textcite{Sellke2024}, for Langevin dynamics of the pure spherical $p$-spin on time scales that do not grow with the dimension, in the limit of zero temperature. The exceptions are instructive. The transition of spiked covariance matrices was found by physicists \parencite{HoyleRattray2004} and by mathematicians \parencite{BaikBenArousPeche2005} within three days of each other, at the turn of February and March 2004 (\cref{ch:spiked}); the theory of the logistic MLE was derived and proved by statisticians, with conic geometry, message passing and the leave-one-out analysis of robust regression, and no replica computation preceded it (\cref{ch:classification}).

The second feature is where the rows that are not proved sit. They cluster in four places: below the de~Almeida--Thouless line and its analogues, where replica symmetry breaks (\cref{ch:sk,ch:gordon}); in the hard phases of inference, where the question is not what can be learned but what can be learned efficiently (\cref{ch:spiked,ch:glm}); away from Gaussian data, where universality fails or is not proved (\cref{ch:lasso,ch:kernels}); and in Part~IV, where the question is where an algorithm goes after a time that grows with the dimension (\cref{ch:kacrice,ch:dmft,ch:sgd}).

\section*{The stone, read across}
\addcontentsline{toc}{section}{The stone, read across}

Each chapter's stone paired the steps of one derivation with the steps of one proof. The same pairs recur from chapter to chapter, and the stone below collects them.

\needspace{12\baselineskip}
\begin{stone}{the physical derivation}{the rigorous route}
\stonestep{Replicate}{Compute $\E Z^n$ for integer $n$: $n$ copies of the system, coupled by the disorder, which interact only through their overlaps $Q_{ab}$ (\cref{ch:freeentropy,ch:rem,ch:sk,ch:gordon}).}{Interpolate: a path from the model to a decoupled system whose derivative is an overlap fluctuation (\cref{ch:sk,ch:nishimori,ch:glm}); or compare: Gordon's inequality replaces the Gaussian matrix by two vectors (\cref{ch:gordon,ch:ridge,ch:lasso,ch:classification}). At $n=2$ the replica computation is the second moment method (\cref{ch:rem}).}
\stonestep{Assume replica symmetry}{All off-diagonal overlaps equal one number $q$, the overlap of two typical configurations.}{A reason for the overlap to concentrate: high temperature (\cref{ch:sk}), the Nishimori identity with a small perturbation (\cref{ch:nishimori,ch:glm}), convexity (\cref{ch:gordon,ch:ridge,ch:lasso,ch:classification}), or a fixed time horizon (\cref{ch:amp,ch:dmft,ch:sgd}). Where there is none, the ansatz fails.}
\stonestep{Find the saddle point}{Extremize the replicated action in $q$; choose the extremum by physical arguments, often a minimum where Laplace's method would ask for a maximum.}{Matching bounds: an inequality with a remainder of fixed sign for every $q$, such as Guerra's bound, and a converse (Aizenman--Sims--Starr, adaptive interpolation, the reverse direction of the CGMT). The theorem decides between minimum and maximum.}
\stonestep{Break the symmetry}{Parisi's hierarchy of overlaps; one step in the REM, with $m=\beta_c/\beta$; a continuum in SK; ultrametricity.}{Poisson--Dirichlet weights and Ruelle's cascades (\cref{ch:rem}); Guerra's broken bound and Talagrand's converse (\cref{ch:sk}); the Ghirlanda--Guerra identities and Panchenko's theorem.}
\stonestep{Continue to $n\to0$}{Take the result for integer $n$, continue it analytically, and divide by $n$.}{No counterpart in this book: no proof continues in $n$. In the REM the exact moment function is not analytic in $n$ (\cref{ch:rem}).}
\stonestep{Remove one variable}{The cavity: the field on a removed spin or sample is Gaussian once the Onsager reaction is subtracted (\cref{ch:amp,ch:lasso,ch:classification}).}{Gaussian conditioning on the past of an iteration (\cref{ch:amp,ch:dmft}); leave-one-out identities for resolvents (\cref{ch:ridge,ch:kernels}).}
\stonestep{Count the states}{The complexity of TAP solutions, or a replicated Kac--Rice computation (\cref{ch:kacrice}).}{The Kac--Rice formula with the GOE determinant for the mean; Subag's second moment for the typical count (\cref{ch:kacrice}).}
\stonestep{Follow the dynamics}{A generating functional, or a dynamical cavity, closes equations for the correlation and the response (\cref{ch:dmft}); for online learning, ordinary differential equations for the order parameters (\cref{ch:sgd}).}{Large deviations of the empirical process; conditioning along an iteration; martingales and Grönwall's lemma (\cref{ch:dmft,ch:sgd}). All on fixed time horizons.}
\end{stone}

Read down the right-hand column, and three things are visible. First, the proofs rarely follow a derivation line by line, but they often follow its architecture. Guerra's interpolation, Gordon's comparison and Bolthausen's conditioning each take one step of the physics, replication, the Gaussian volume integral and the cavity, and turn it into an identity or an inequality whose remainder has a sign. Second, the replica-symmetric ansatz never becomes a lemma. Every proof of a replica-symmetric formula in this book replaces the ansatz by a reason for the overlap to concentrate, and the book met four reasons: high temperature, Bayes optimality, convexity, and a fixed horizon. Where none of them holds, as below the de~Almeida--Thouless line, for the perceptron with a negative margin, and in the long-time dynamics of rugged landscapes, either replica symmetry is false or nothing is proved. Third, the continuation $n\to0$, the step that mathematicians distrust most, is the one step that no proof reproduces. The rigorous theory does not justify the replica trick; it makes the trick unnecessary, one model at a time.

\section*{What the proofs assume}
\addcontentsline{toc}{section}{What the proofs assume}

The twenty-four rows marked \emph{proved under} name the hypotheses that the theorems need and the predictions do not, and the same few hypotheses recur.

\emph{Gaussian data.} Gordon's comparison is a statement about Gaussian processes, and the main theorems of Part~III start from Gaussian inputs. Universality theorems extend some of them to independent entries (\cref{ch:ridge,ch:lasso}) and to random features (\cref{ch:kernels}), but not to general isotropic designs, for which the Gaussian risk of least squares provably fails and that of the LASSO fails in simulations \parencite{HanShen2023} (\cref{ch:lasso}), and not to real data.

\emph{Bayes optimality.} The proofs of the replica formulas of Part~II use the Nishimori identity, which holds only when the statistician knows the prior and the channel. At the wrong temperature the replica-symmetric formula is false (\cref{ch:nishimori}), and the mismatched student of \cref{ch:glm} is open.

\emph{Convexity.} The two-sided comparison of \cref{ch:gordon} needs convex sets and convex--concave objectives. Without convexity Gordon gives one side, which is why Gardner's formula for the negative perceptron is only an upper bound on its capacity.

\emph{Fixed horizons.} State evolution holds for a fixed number of iterations as the dimension grows (\cref{ch:amp}), dynamical mean-field theory for fixed times, at positive temperature or with smooth losses (\cref{ch:dmft}), and the Saad--Solla equations on fixed intervals of rescaled time (\cref{ch:sgd}). The physically interesting questions, such as where gradient flow ends, how long a plateau lasts, or whether AMP converges, are about the opposite order of limits.

\emph{Numerical evaluation.} Several thresholds are theorems up to the numerical solution of a scalar equation: the information-theoretic thresholds of sparse PCA (\cref{ch:spiked}), $\alpha_{\mathrm{IT}}=1.2447$ for the Ising teacher (\cref{ch:glm}), and the capacity $0.8331$ of the Ising perceptron, whose proof rests on numerical conditions (\cref{ch:gordon}). In each case the book's scripts recompute the number; in \cref{ch:glm} the recomputation disagrees with the value printed in the source.

\section*{Where the physics was wrong}
\addcontentsline{toc}{section}{Where the physics was wrong}

Nineteen rows are marked \emph{known false}. They are of three kinds, and only the last is a failure of physics, or of anyone, in the ordinary sense.

Seven rows record replica-symmetric answers that physics itself rejects. Below the critical temperature of the REM the replica-symmetric entropy is negative (\cref{ch:rem}); below the de~Almeida--Thouless line of SK the replica-symmetric saddle point is unstable (\cref{ch:sk}); for the perceptron with a negative margin, and for the fewest errors beyond capacity, the replica-symmetric solution is unstable (three rows: the prologue, \cref{ch:gordon,ch:classification}); for binary weights the replica-symmetric capacity $4/\pi$ exceeds the counting bound $1$ (\cref{ch:gordon}); for $\ell_0$ reconstruction the replica-symmetric success solution is unstable at every load (\cref{ch:lasso}). In each case physicists found the failure themselves. For the REM, SK and binary weights they also found the broken solution; for the negative perceptron they predict full replica-symmetry breaking, and for $\ell_0$ reconstruction the broken computation was announced and is still open. Theorems confirm the failure for the REM, for SK, for binary weights and for the negative perceptron, the last so far only below $\kappa=-0.622$ and as $\kappa\to-\infty$; for the fewest errors beyond capacity and for $\ell_0$ reconstruction the evidence is the physicists' own stability computation. Where theorems reach the broken solution at all, they confirm it too.

Four record shortcuts whose failure the physics itself explains: the annealed average below the critical temperature (\cref{ch:freeentropy}) and in learning curves (\cref{ch:glm}), naive mean field without the Onsager term (\cref{ch:amp}), and the naive continuation of the integer moments of the REM to $n\to0$ (\cref{ch:rem}).

The remaining eight record extrapolations beyond the regime in which a derivation or a theorem holds, made by physicists and mathematicians alike: the Gaussian approximation to learning curves at zero ridge in fixed dimension, which fails on the circle, where the interpolant is nearly piecewise linear and its noise term is the $2\sigma^2/3$ of linear interpolation \parencite{BelkinEtAl2018}, not the formula's $\sigma^2$ (\cref{ch:kernels}); linear Gaussian equivalence in polynomial scalings or after one large gradient step, and the Gaussian covariate model for the classification error of real feature maps on real images (\cref{ch:kernels}); universality of the Gaussian formulas for every isotropic design, which fails for least squares and, in simulations, for the LASSO (\cref{ch:lasso}); equality of the typical and mean numbers of critical points in mixed models (\cref{ch:kacrice}); equality of the long-time dynamics and the statics (\cref{ch:dmft}); and the claim that $d^{k-1}$ samples are necessary for every gradient method to learn a single-index target of information exponent $k$, which smoothing and the reuse of samples both beat (\cref{ch:sgd}). The replica-symmetric formula at the wrong temperature (\cref{ch:nishimori}) belongs here too: the formula is right on the Nishimori line and was never claimed off it.

\section*{The preprints}
\addcontentsline{toc}{section}{The preprints}

Nine rows rest on unrefereed preprints. Four of them are preprints posted between 2013 and 2024 whose results, as the boards use them, have no refereed version: El~Karoui's analysis of unregularized robust regression with non-Gaussian designs, whose published revision keeps only ridge-regularized estimators (\cref{ch:lasso}), the dynamical mean-field theory of gradient flow of \textcite{CelentanoChengMontanari2021} (\cref{ch:dmft}), the long-horizon AMP of \textcite{LiWei2022} (\cref{ch:amp}), and the deterministic equivalent for general kernels of \textcite{MisiakiewiczSaeed2024} (\cref{ch:kernels}). The other five appeared in 2026 and settle questions, most of them open for decades: replica symmetry of SK on the whole high-temperature side of the de~Almeida--Thouless line with a field \parencite{Lopatto2026a,KusuokaNakajima2026b}; full replica-symmetry breaking below it \parencite{Lopatto2026b,Lopatto2026c}; the convergence of the overlap of SK itself to the Parisi measure \parencite{AronowLopatto2026b}; large deviations of the SK free entropy at the order predicted by \textcite{ParisiRizzo2009} \parencite{AronowLopatto2026a,Pei2026}; and the capacity of the Ising perceptron without numerical conditions \parencite{Shmalo2026}. \Cref{tab:epilogue:preprints} lists them with their arXiv versions as of 25 October 2026, together with one more 2026 preprint that \cref{ch:sk} reports in its text but not on its board.

\begin{table}[t]
\begin{fullwidth}
\footnotesize
\renewcommand{\arraystretch}{1.2}
\begin{tabularx}{\linewidth}{@{}>{\raggedright\arraybackslash}p{38mm}>{\raggedright\arraybackslash}p{30mm}>{\raggedright\arraybackslash}X@{}}
\toprule
\textsc{\color{stone}result} & \textsc{\color{stone}preprint} & \textsc{\color{stone}what the authors say about AI assistance}\\
\midrule
Replica symmetry of SK for $h>0$ wherever $\beta^2\E\,\mathrm{sech}^4\le1$ (\cref{ch:sk}) & \textcite{Lopatto2026a}, arXiv:2604.11921v2 & written with the assistance of large language models, which suggested arguments, helped draft and revise, and coded simulations\\
The same without the Parisi formula, with rate $1/N$ where $\beta^2\E\,\mathrm{sech}^4<1$ & \textcite{KusuokaNakajima2026b}, arXiv:2608.23413v2 & for their key step the authors \enquote{asked ChatGPT 5.6 Sol to construct a proof}; they revised the argument in detail, take full responsibility for it, and formalized the proof in Lean~4 with Codex\\
Full RSB at $h=0$ for every $\beta>1$ & \textcite{Lopatto2026b}, arXiv:2607.11756v4 & written with the assistance of large language models\\
The Parisi measure for $h>0$ has an interval as support, with a density inside & \textcite{Lopatto2026c}, arXiv:2610.00127v1 & written with the assistance of large language models; the observables of the proof were located with \enquote{a massive number of numerical experiments orchestrated by large language models}\\
The zero-temperature Parisi measure has support $[0,1)$ (in the text of \cref{ch:sk}) & \textcite{ChenHB2026}, arXiv:2607.18032v1 & the arXiv record states that \enquote{ChatGPT 5.6 is the actual author}\\
The overlap law of SK is the Parisi measure & \textcite{AronowLopatto2026b}, arXiv:2609.17442v1 & written with the assistance of large language models\\
Upper tail $e^{-Na\delta^{6/5}}$ of the SK free entropy (\cref{ch:freeentropy}) & \textcite{AronowLopatto2026a}, arXiv:2609.09103v1; \textcite{Pei2026}, arXiv:2610.06480v2 & Aronow and Lopatto: written with the assistance of large language models; Pei: none stated\\
Capacity $0.8330786$ of the Ising perceptron (\cref{ch:gordon}) & \textcite{Shmalo2026}, arXiv:2607.11010v1 & \enquote{Most of the work reported here was done by two AI systems}; the proof has not yet been independently reproduced or peer reviewed\\
\bottomrule
\end{tabularx}
\end{fullwidth}
\caption{The 2026 preprints that the chapters report, with the arXiv versions current on 25 October 2026 and what each says about the use of AI systems, quoted or paraphrased from the papers as the chapters report them. None is refereed. The boards mark every result in the table \emph{preprint} except the fifth, which \cref{ch:sk} reports only in its text.}\label{tab:epilogue:preprints}
\end{table}

These results are reported as their authors state them, and the boards do not count them as proved. By their authors' accounts, AI systems also did a substantial part of the mathematical work in several of them. The book does not treat them as more or less likely to be right than other preprints on that account. What it changes is the weight of independent checking, and two of them invite such checking in forms that are unusual in the field: a formalization in Lean (\cref{ch:sk}) and a computer-assisted proof (\cref{ch:gordon}).

\section*{Open problems}
\addcontentsline{toc}{section}{Open problems}

Each chapter closed with its open problems. Ten of them, in the chapters' order, show the shape of what is not known as of October 2026.
\begin{enumerate}
\item A mathematical framework for the replica method with broken symmetry, which would explain, among other things, why the extremum in the REM is a minimum over $[n,1]$ \parencite[notes to ch.~8]{MezardMontanari2009}; and, more modestly, a class of disordered systems in which the moment problem for $Z_N$ is solved, so that the replica trick is justified \parencite{Tanaka2007} (\cref{ch:freeentropy,ch:rem}).
\item The order of the fluctuations of $\log Z_N$ for SK below the critical temperature, whose variance the preprint of \textcite{AronowLopatto2026a} places between the orders $N^{4/15}$ and $N^{7/15}$ (\cref{ch:freeentropy,ch:sk}).
\item Whether the local minimizer of the TAP free energy near the posterior mean of $\Z_2$ synchronization, which \textcite{CelentanoFanMei2023} construct for every $\lambda>1$, is global, as they conjecture from numerics (\cref{ch:nishimori,ch:amp}).
\item Whether some polynomial-time algorithm beats message passing in the hard phases of sparse PCA \parencite{LesieurKrzakalaZdeborova2017} and of the Ising teacher--student perceptron \parencite{BarbierEtAl2019} (\cref{ch:spiked,ch:glm}).
\item The capacity of the spherical perceptron with a negative margin, for which physics predicts full replica-symmetry breaking \parencite{FranzEtAl2017} and rigorous bounds match only to leading order as $\kappa\to-\infty$ \parencite{MontanariZhongZhou2024} (\cref{ch:gordon}).
\item Universality of the existence threshold of the logistic MLE for non-Gaussian inputs \parencite[Sec.~5]{CandesSur2020}, and a characterization of the designs that share the Gaussian LASSO risk (\cref{ch:lasso,ch:classification}).
\item A criterion, for a given kernel and data distribution, for the exactness of the Gaussian approximation to the learning curve, and the exchange of the limits $\lambda\to0$ and $d\to\infty$ for random features \parencite[Remark~1]{MeiMontanari2022} (\cref{ch:kernels}).
\item The typical complexity of the spiked tensor model near its isolated minimum, and in its band of minima near the signal strength at which that minimum appears, which \textcite{RosEtAl2019} computed by a replicated Kac--Rice formula and \textcite{BenArousMeiMontanariNica2019} leave open (\cref{ch:kacrice}).
\item That gradient flow ends at the threshold energy of the pure spherical $p$-spin, which \textcite{Sellke2024} proves only for Langevin dynamics on time scales that do not grow with the dimension, in the limit of zero temperature, and not for the flow itself (his Remark~1.5); and the dynamical transition and aging for $p\ge3$ (\cref{ch:kacrice,ch:dmft}).
\item The length of the plateau of online SGD on a committee machine, of order $\log d$ and with a limiting law, from a random start \parencite[Section~3]{VeigaEtAl2022} (\cref{ch:sgd}).
\end{enumerate}
Many of these share one feature: they are about the regime in which the four reasons for concentration of the previous sections fail, or about times that grow with the dimension. The physics of these regimes is rich and, in most cases, believed; the mathematics is where the frontier is.

\section*{What to read next}
\addcontentsline{toc}{section}{What to read next}

On the physical side, \textcite{MezardParisiVirasoro1987} collect the founding papers of spin-glass theory with commentary, \textcite{Nishimori2001} develops the Nishimori line, \textcite{EngelVandenBroeck2001} is the textbook of the statistical mechanics of learning, \textcite{MezardMontanari2009} treat the cavity method and message passing for inference and coding, and \textcite{ZdeborovaKrzakala2016} review the statistical physics of inference. On the mathematical side, \textcite{Talagrand2011a,Talagrand2011b} and \textcite{Panchenko2013b} develop the theory of mean-field spin glasses, and \textcite{Bovier2006} that of the random energy models. Between them, \textcite{MontanariSen2024} teach replica and cavity computations alongside interpolation bounds, on two models, for readers without physics, the lecture notes of \textcite{KrzakalaZdeborova2024} follow several replica computations with their rigorous counterparts, \textcite{BandeiraPerryWein2018} explain the physicists' heuristics for computational thresholds to mathematicians, \textcite{FengEtAl2022} unify the theory of approximate message passing, and \textcite{ThrampoulidisAbbasiHassibi2018} set out the convex Gaussian min--max theorem with its applications to estimation. For dynamics, \textcite{Cugliandolo2024} surveys dynamical mean-field methods and their recent applications.

\printchapterbib
\end{refsection}
\rstfloatrestore
